% Dependency-free 16:9 slide deck (one landscape page per slide).
% Uses only graphicx / geometry / xcolor — no beamer, no package installs.
\documentclass[11pt]{article}
\usepackage[T1]{fontenc}
\usepackage[paperwidth=25.6cm,paperheight=14.4cm,margin=1.1cm]{geometry} % 16:9
\usepackage{graphicx}
\usepackage{xcolor}
\usepackage[hidelinks]{hyperref}
\renewcommand{\familydefault}{\sfdefault}
\pagestyle{empty}
\setlength{\parindent}{0pt}
\graphicspath{{../paper_figures/}}

\definecolor{teal}{RGB}{14,140,127}
\definecolor{ctrl}{RGB}{21,101,192}
\definecolor{adred}{RGB}{198,40,40}
\definecolor{ink}{RGB}{22,32,47}
\color{ink}

\newcounter{sld}
\newcommand{\pageno}{\par\vspace{0.15cm}\hfill{\scriptsize\color{gray!70}\thesld}}
\newcommand{\eyebrow}[1]{{\footnotesize\color{teal}\textbf{\uppercase{#1}}}\par\vspace{0.2cm}}

% a figure slide: title, figure, one-line caption
\newcommand{\figslide}[3]{%
  \stepcounter{sld}%
  \noindent{\large\color{teal}\textbf{#1}}\par\vspace{0.25cm}%
  \begin{center}\includegraphics[width=0.82\linewidth,height=0.66\paperheight,keepaspectratio]{#2}\end{center}%
  \vfill{\small #3}\par\pageno\clearpage}

\begin{document}

% 1 TITLE
\stepcounter{sld}
\vspace*{2.2cm}
\eyebrow{Personalised models of resting-state fMRI \textperiodcentered\ Alzheimer's disease}
{\fontsize{30}{34}\selectfont\bfseries Optimal stimulation sites\\ are not the most affected}\par
\vspace{1.4cm}
{\small\color{ink!85} C.~Capone \quad E.~Cece \quad A.~Ciardiello \quad G.~Gigante \quad E.~Cisbani \quad M.~Mattia}\par
\vspace{0.15cm}{\small\color{gray} Istituto Superiore di Sanit\`a, Rome}
\clearpage

% 2 PROBLEM
\stepcounter{sld}
\vspace*{2.6cm}
\begin{center}
{\Large Alzheimer's is a \textcolor{adred}{distributed} network disorder.}\\[0.5cm]
{\Large Yet targeted neuromodulation needs a \textbf{focal} target.}\\[1.1cm]
{\fontsize{26}{30}\selectfont\bfseries\color{teal} Where should we stimulate?}
\end{center}
\pageno\clearpage

% 3-8 FIGURE SLIDES
\figslide{Resting-state fMRI: controls and patients}{figure1_data}%
  {Individual resting-state fMRI from cognitively unimpaired controls and AD patients.}

\figslide{A perturbable digital twin of each patient}{figure2_model}%
  {A personalised reservoir model whose free-running dynamics \textbf{reproduce the patient's own} resting-state activity.}

\figslide{Reading the disease from the model}{figure3_classification}%
  {The fitted model parameters separate AD from controls --- a \textit{target} for intervention, not a diagnostic end.}

\figslide{Correcting the connectivity reverts the signature}{figure4_stimulation}%
  {Shifting the whole connectivity toward controls works; \textcolor{adred}{focal edits fail at any amplitude} --- the correction is distributed.}

\figslide{The best target is not the most affected site}{figure5_topsites}%
  {\textcolor{adred}{Most-changed} connectivity sites and \textcolor{teal}{effective} stimulation targets are largely disjoint.}

\figslide{One site, and closed-loop control}{figure6_singlecompare}%
  {A resonant drive at each patient's discriminant-aligned site reverts the classification; closed-loop titration reaches it at lower dose.}

% 9 TAKEAWAY
\stepcounter{sld}
\vspace*{2.6cm}
\begin{center}
{\fontsize{24}{30}\selectfont\bfseries Stimulate where the network \textcolor{teal}{responds} ---}\\[0.3cm]
{\fontsize{24}{30}\selectfont\bfseries not where it is most \textcolor{adred}{affected}.}\\[1.2cm]
{\normalsize\color{ink!85} Model-informed, personalised targeting turns ``where to stimulate'' into a computable prescription.}
\end{center}
\pageno\clearpage

\end{document}
